\documentclass[11pt]{article}
\usepackage[a4paper,margin=29mm]{geometry}
\usepackage{amsmath,amssymb,amsthm,mathtools}
\usepackage{microtype}
\usepackage[hidelinks]{hyperref}
\setlength{\emergencystretch}{2em}

\newtheorem{theorem}{Theorem}[section]
\newtheorem{proposition}[theorem]{Proposition}
\newtheorem{lemma}[theorem]{Lemma}
\newtheorem{corollary}[theorem]{Corollary}
\theoremstyle{definition}
\newtheorem{definition}[theorem]{Definition}
\newtheorem{example}[theorem]{Example}
\theoremstyle{remark}
\newtheorem{remark}[theorem]{Remark}
\newcommand{\R}{\mathbb R}
\newcommand{\Q}{\mathbb Q}
\newcommand{\Z}{\mathbb Z}
\newcommand{\one}{\mathbf 1}
\newcommand{\Ch}{\mathcal C}
\newcommand{\NR}{N_{\R}}
\newcommand{\MR}{M_{\R}}
\newcommand{\pos}{\operatorname{pos}}
\newcommand{\Span}{\operatorname{span}_{\R}}
\newcommand{\aff}{\operatorname{aff}}
\newcommand{\lin}{\operatorname{lin}}
\newcommand{\supp}{\operatorname{supp}}
\newcommand{\ind}{\mathbf 1}

\title{Weighted tropical Chow reconstruction for space curves}
\author{YINGFENG JIANG}
\date{October 2, 2026}

\begin{document}
\maketitle

\begin{abstract}
We prove that the weighted tropical Chow transform determines a tropical
curve in $\R^4/\R\one$, without requiring its support to be a fan.
More generally, the transform is injective on finite rational polyhedral
one-chains with signed rational weights, whether or not they are balanced.
The proof separates contributions according to their affine
three-dimensional spans. For edges outside the coordinate directions,
small parameters in a suitable incidence cone localize every competing
source point. For coordinate directions, the remaining contributions
split into a term independent of a thickness parameter and a thin band.
The thin band tends to zero locally in $L^1$, while its normalized measure
converges to the source length measure. This eliminates the remaining
kernel without any general-position or realizability assumption.
\end{abstract}

\section{Introduction}

The tropical Chow construction associates to a tropical variety a weighted
incidence cycle in a tropical Grassmannian. For a curve
$\Sigma\subset\R^4/\R\one$, this cycle parametrizes tropical lines meeting
$\Sigma$, with multiplicities determined by a weighted Minkowski sum.
In the public arXiv version~2 of his paper, Tripoli formulated the
injectivity problem and proved the space-curve case when the support is
a fan \cite[Conjecture~3.4 and Theorem~3.9]{Tripoli}.
All numbered references to that paper below use this same public version.

Our purpose is to remove the fan hypothesis for curves with finitely many
polyhedral cells. All equalities below retain weights and are taken after
common subdivision. In particular, positions in the ambient space are
part of the data.

\begin{theorem}\label{thm:curves}
Let $\Sigma$ and $\Sigma'$ be finite rational polyhedral tropical curves
in $\R^4/\R\one$, with positive integral weights. Then
\[
 Z_\Sigma=Z_{\Sigma'}\quad\Longrightarrow\quad\Sigma=\Sigma'.
\]
The curves may have bounded edges, distinct vertices, and more than one
connected component.
\end{theorem}

Finiteness means a finite polyhedral structure, not compact support.
Thus rays and complete affine lines are allowed. Rational polyhedra are
taken with respect to the indicated quotient lattices. We work in the
finite-coordinate tropical tori, without additional strata at infinity.

The proof uses a linear extension of the Chow transform to signed
polyhedral chains. Such chains are unoriented weighted polyhedra modulo
subdivision; balancing is not imposed. Write $\Ch_d(V)$ for their
$\Q$-vector space in a real vector space $V$ with a fixed lattice.
Section~\ref{sec:transform} defines a map
\[
 T:\Ch_1(\NR)\longrightarrow\Ch_3(\MR),\qquad
 N=\Z^4/\Z\one,\quad M=\Z^6/\Z\one.
\]

\begin{theorem}\label{thm:chains}
The map $T$ is injective.
\end{theorem}

This stronger statement is useful because separating contributions in the
target does not preserve balancing. We therefore carry out that separation
in the vector space of chains, and recover Theorem~\ref{thm:curves} by
restricting to balanced inputs.

In the proof of Lemma~3.8, Tripoli already extends the transform to
rationally weighted, possibly unbalanced fans supported on the eight rays
$\R_{\ge0}e_i$ and $\R_{\ge0}(-e_i)$, $i=0,\ldots,3$, all based at the
origin, and proves its injectivity \cite{Tripoli}. Our extension allows
arbitrary finite affine one-chains; thin-band recovery handles coordinate
edges in arbitrary positions.

Although every tropical cycle in a vector space can be expressed as a
sum of translated fan cycles \cite[Theorem~5.4]{AllermannHampeRau}, this does
not by itself reduce the question to injectivity on fans. Images of distinct
translated summands may cancel. Our argument controls precisely these
overlaps; it does not require choosing a decomposition into fan cycles.

There are two parts to the proof. First, different affine
three-dimensional spans in the target can meet only in lower dimension.
Within a suitable span, the contributions of an edge not parallel to a
coordinate direction come from just one incidence cone. A short-parameter
argument then rules out cancellation by distant source points.
Second, for coordinate directions, several cones do contribute to the
same span. Their weight function has the form
\[
 K_t=H+C_t,
\]
where $H$ contains all side-cone contributions and does not depend on $t$.
The central contribution $C_t$ is a thin band swept out by translating
the source through a segment of length $t$. A boundary limit first removes
$H$ and then recovers the source measure.

Weights are essential: the injectivity statement for incidence supports
alone already fails for fan curves
\cite[Example~3.5]{Tripoli}. The result concerns the full weighted Chow
cycle, rather than a Chow polytope or its normal complex; see
\cite{Fink} for the latter construction. Our argument establishes
uniqueness. It does not assert an algorithm, or a complexity bound, for
reconstructing an arbitrary input from its output.

\section{The chain transform and the incidence fan}\label{sec:transform}

\subsection{Polyhedral chains and weights}

For a rational affine polyhedron $P$ in a lattice vector space $V$ with
lattice $L$, put
\[
 \lin(P)=\Span(P-P),\qquad L_P=L\cap\lin(P).
\]
In particular, the lattice of an affine cell is its direction lattice,
not the group generated by its affine integral points.

\begin{definition}
The vector space $\Ch_d(V)$ consists of finite rational linear
combinations of symbols $[P]$ for rational $d$-dimensional polyhedra,
modulo subdivision: if $P$ is subdivided into $d$-dimensional cells $P_i$,
then $[P]=\sum_i[P_i]$. Cells of dimension less than $d$ are discarded.
Equivalently, an element is a finite polyhedral complex with signed
rational weights on its $d$-cells, modulo common subdivision and removal
of zero-weight cells. No orientation or balancing condition is included.
\end{definition}

For a chosen representative of a one-chain $D$, write $D^{(1)}$ for its
one-dimensional cells. Such a chain is balanced if, at every vertex $x$,
\[
 \sum_{e\ni x}w_e u_{e,x}=0,
\]
where $u_{e,x}$ is the primitive direction of $e$ pointing away from $x$.
One may subdivide complete lines to state this condition. A tropical
curve in this paper is an effective balanced one-chain with integral
weights; its vertices carry no additional zero-dimensional weights.
The support is the closure of the nonzero-weight one-cell interiors.
At a point in the relative interior of a one-cell, a sufficiently small
neighborhood of this support consists of that cell alone.

Let $e_i$ denote the images of the standard basis of $\Z^4$ in $N$.
Use the order $01,02,03,12,13,23$ for the coordinates of $M$, and write
$E_{ij}$ for its standard generators. Pair subscripts are unordered,
so $E_{ij}=E_{ji}$ and the same convention applies to pair coordinates.
Define
\[
 \varphi:N\longrightarrow M,\qquad
 [u_0,u_1,u_2,u_3]\longmapsto[u_i+u_j]_{i<j},
 \qquad f_i=\varphi(e_i).
\]
The map is well defined. It is injective and has saturated image: for
distinct $i,j$ and a third index $k$,
\[
 u_i-u_j=\varphi(u)_{ik}-\varphi(u)_{jk}.
\]
Consequently, if $\varphi(u)$ represents an integral point of $M$, all
differences $u_i-u_j$ are integral, so $u$ represents a point of $N$.

Let $\Gamma_0$ be the tropical incidence fan of lines through the origin,
with its usual multiplicities. We use the minimum convention for
tropicalization; in this convention the standard tropical line has rays
in directions $e_0,e_1,e_2,e_3$. The constant-coefficient incidence space
may be taken over $\mathbb C$, and rational valuations realized over
the Puiseux series field. Its explicit incidence cones are given below.
For a one-cell $e$ and a two-cone $\tau$ of $\Gamma_0$, set
\[
 P_{e,\tau}=\varphi(e)+\tau.
\]
If this polyhedron has dimension three, define
\begin{equation}\label{eq:index}
 m(e,\tau)=
 \bigl[M\cap\lin(P_{e,\tau}):\varphi(N_e)+M_\tau\bigr].
\end{equation}
This is a positive finite integer. For a chain $D=\sum_e w_e[e]$, define
\begin{equation}\label{eq:transform}
 T(D)=\sum_{\substack{e\in D^{(1)},\ \tau\in\Gamma_0^{(2)}\\
                         \dim P_{e,\tau}=3}}
 w_e w_\tau m(e,\tau)[P_{e,\tau}],
 \qquad w_\tau=1.
\end{equation}
Overlaps in this sum are added after common subdivision.

\begin{lemma}\label{lem:well-defined}
Formula~\eqref{eq:transform} defines a linear map on $\Ch_1(\NR)$
and is unchanged by subdivision of the incidence fan. For a balanced
input it agrees with the weighted Minkowski transform
$Z_D=\varphi(D)+\Gamma_0$.
\end{lemma}

\begin{proof}
On a product $e\times\tau$ with three-dimensional image, the addition
map $(x,g)\mapsto\varphi(x)+g$ is injective on the direction space:
both domain and image have dimension three. A subdivision of the product
therefore subdivides its image, with the same direction lattices and the
same index~\eqref{eq:index}. Images of subdivision boundaries have
dimension at most two. If the original image has smaller dimension,
every subdivided image does too, and all are discarded. This proves
well-definedness and subdivision invariance; linearity follows from the
weight formula. On balanced chains the formula is the standard
full-rank weighted Minkowski sum, as in
\cite[Section~2, equation~(5) and Definition~2.7]{Tripoli}.
\end{proof}

\begin{remark}
For a nonbalanced chain we use the explicit linear map
\eqref{eq:transform}, not a claim that the input is a tropical cycle.
There is no additional division by a covering degree in this definition.
\end{remark}

\subsection{The fifteen cones}

For a partition $\{a,b\}\sqcup\{c,d\}=\{0,1,2,3\}$, write
\[
 h_{ab}=e_a+e_b,
 \qquad H_{ab}=\Span(e_a,e_b).
\]
An identity used throughout is
\begin{equation}\label{eq:identity}
 \varphi(h_{ab})=f_a+f_b=E_{ab}-E_{cd}.
\end{equation}

\begin{lemma}\label{lem:fan}
The fan $\Gamma_0$ admits a structure with the following fifteen maximal
cones, each of dimension two and weight one:
\begin{align}
 \sigma_{ab|cd}&=\pos(E_{ab},E_{cd})
       &&\text{for the three unordered partitions }ab|cd,
       \label{eq:central}\\
 \tau_{ij;k}&=\pos(E_{ij},-f_k)
       &&\text{for }k\notin\{i,j\}.
       \label{eq:side}
\end{align}
In each case the displayed generators form a basis of the saturated
direction lattice of the cone.
\end{lemma}

\begin{proof}
The incidence linear space has Pl\"ucker coordinates
$p_{ij}=z_j-z_i$: a line through $[1:1:1:1]$ is spanned by that point
and $[z_0:z_1:z_2:z_3]$. Set $w_{ij}=\operatorname{val}(z_j-z_i)$
and normalize $\min_{i<j}w_{ij}=0$. For any three labels, the minimum
of their three pair valuations is attained at least twice. Thus for
each $r\ge0$, the relation $i\sim_r j$ defined by $w_{ij}>r$, with
$i\sim_r i$, is an equivalence relation. Its nontrivial classes form
nested proper clusters as $r$ increases.

On four labels there are at most two nontrivial proper clusters in
a maximal such hierarchy. They are either two disjoint two-element
clusters, or a three-element cluster with a two-element subcluster.
Indeed, two proper clusters which intersect are nested, a proper cluster
has at most three labels, and a three-element cluster contains at most
one two-element cluster in a nested hierarchy. The first pattern
gives~\eqref{eq:central}.
For the second, the valuation vector of the three-element cluster
complementary to $k$ is
\[
 \sum_{i<j,\ i,j\ne k}E_{ij}=-f_k,
\]
and its subcluster contributes $E_{ij}$, giving~\eqref{eq:side}.
Hierarchies with fewer clusters give faces. Conversely, the interiors
of the two representative cones are realized, for rational
$\alpha,\beta,r,s>0$, by
\[
 (z_0,z_1,z_2,z_3)=(0,t^\alpha,1,1+t^\beta),\qquad
 (z_0,z_1,z_2,z_3)=(0,t^{r+s},1,t^s),
\]
respectively. These valuation vectors are
$\alpha E_{01}+\beta E_{23}$ and $rE_{01}-sf_2$ in $\MR$.
Permuting the labels and taking closures gives all fifteen cones.

For completeness, their multiplicities can be read from two initial
ideals. The incidence ideal is generated by the three independent
linear relations
\[
 p_{12}-p_{02}+p_{01},\qquad
 p_{13}-p_{03}+p_{01},\qquad
 p_{23}-p_{03}+p_{02}.
\]
At an interior point of $\sigma_{01|23}$ its initial ideal is
\[
 (p_{12}-p_{02},\ p_{13}-p_{03},\ p_{02}-p_{03}),
\]
and at an interior point of $\tau_{01;2}$ it is
\[
 (p_{12}-p_{02},\ p_{13}-p_{03},\ p_{23}+p_{02}).
\]
The displayed initial forms lie in the respective initial ideals.
In each case they generate an ideal of three independent linear forms,
with the same ordinary graded Hilbert series as the incidence ideal.
Preservation of that Hilbert series under taking weight initial ideals
shows that these are the entire initial ideals. Both define reduced
irreducible linear spaces with nonempty coordinate-torus intersections,
so their torus degenerations have multiplicity one. Permutations of
labels give the same conclusion for every maximal cone.
The fan description, including these weights, is also recorded in
\cite[Section~3.1, Figure~3]{Tripoli}.

For the lattice assertion it suffices, by permutations of the four
labels, to check $(E_{01},E_{23})$ and $(E_{01},-f_2)$.
Identify $M$ with $\Z^5$ by subtracting the $23$ coordinate and then
deleting it. The first pair has a $2\times2$ minor of absolute value
one in rows $01,02$, and the second in rows $01,03$.
Thus each pair generates its saturated rank-two lattice.
\end{proof}

To keep track of the three partitions, consider the quotient map
\begin{equation}\label{eq:quotient}
 \pi:\MR\longrightarrow\R^2,\qquad
 p\longmapsto
 (p_{01}+p_{23}-p_{02}-p_{13},\,
  p_{01}+p_{23}-p_{03}-p_{12}).
\end{equation}
Its kernel is $\varphi(\NR)$: the inclusion follows from the formula
for $\varphi$, and both spaces have dimension three. The images of
$E_{01},E_{02},E_{03}$ are respectively
$(1,1),(-1,0),(0,-1)$, while complementary pairs have equal images.
In particular the three partition directions are pairwise noncollinear.

The partition criterion in the following lemma is equivalent to
\cite[Lemma~3.7]{Tripoli}. Its choice of central cone and separation of
linear spans follow the proof of \cite[Theorem~3.9]{Tripoli}; we give a
self-contained argument in our notation. The passage to arbitrary affine
sources will additionally require grouping by exact affine span
(Lemma~\ref{lem:separation}) and localizing competing source points through
the equality $\alpha'+\beta'=\alpha+\beta$
(Proposition~\ref{prop:ordinary}, equation~\eqref{eq:localization}).

\begin{lemma}\label{lem:ordinary-linear}
Let $v\in\NR$ be nonzero and not parallel to any $e_i$. There is a
partition $ab|cd$ such that $v\notin H_{ab}\cup H_{cd}$. For this
partition the space
\[
 U=\Span(\varphi(v),E_{ab},E_{cd})
\]
has dimension three, contains none of the $f_i$, and contains exactly
one maximal cone of $\Gamma_0$, namely $\sigma_{ab|cd}$.
Moreover,
\begin{equation}\label{eq:intersection}
 U\cap\varphi(\NR)=\varphi(\Span(v,h_{ab})).
\end{equation}
\end{lemma}

\begin{proof}
Membership in $H_{ab}$ means that the complementary two coordinates
of a representative of $v$ are equal. Since $v$ is not a coordinate
direction, no three of its four coordinates are equal. We can therefore
partition the four indices into two pairs within each of which the
coordinates differ. This gives the required partition.

The vectors $v,h_{ab}$ are independent, since $h_{ab}$ belongs to both
$H_{ab}$ and $H_{cd}$. Equations~\eqref{eq:identity}
and~\eqref{eq:quotient} give the dimension statement
and~\eqref{eq:intersection}. If $f_i\in U$, then
$e_i\in\Span(v,h_{ab})$. As $e_i$ is independent of $h_{ab}$, this
would imply $v\in H_{ab}$ for $i\in\{a,b\}$, or $v\in H_{cd}$
otherwise. Hence no side cone is contained in $U$. The quotient
directions in~\eqref{eq:quotient} exclude the other central cones.
\end{proof}

\begin{lemma}\label{lem:coordinate-linear}
For a pair $a,b$ with complementary pair $c,d$, put
\[
 U_{ab}=\Span(f_a,f_b,E_{cd}).
\]
Then $U_{ab}\cap\varphi(\NR)=\varphi(H_{ab})$, and
$f_i\in U_{ab}$ if and only if $i\in\{a,b\}$. The maximal cones
of $\Gamma_0$ contained in $U_{ab}$ are exactly
\begin{equation}\label{eq:three-cones}
 \sigma_{ab|cd},\quad
 \pos(E_{cd},-f_a),\quad\pos(E_{cd},-f_b).
\end{equation}
The vectors $f_a,f_b,E_{cd}$ form a basis of $M\cap U_{ab}$.
Each full-rank pairing of a source direction $e_a$ or $e_b$ with a
cone in~\eqref{eq:three-cones} has index one
in~\eqref{eq:index}.
\end{lemma}

\begin{proof}
The intersection statement follows by applying $\pi$, and the assertion
about the $f_i$ follows since any three of the $e_i$ form a basis of
$N$. The quotient directions restrict possible cones to the partition
$ab|cd$. Among its side cones, exactly those displayed have their
coordinate generator in $U_{ab}$.

By symmetry take $ab=01$, $cd=23$. In the quotient coordinates used in
Lemma~\ref{lem:fan}, the columns $f_0,f_1,E_{23}$ have the minor
\[
 \det\begin{pmatrix}
 1&1&-1\\1&0&-1\\0&1&-1
 \end{pmatrix}=1
\]
in rows $01,02,12$. Thus they form a basis of the saturated lattice.
In the basis $f_a,f_b,E_{cd}$, identity~\eqref{eq:identity} reads
$E_{ab}=(1,1,1)$. Each of the four noncollapsed source--cone pairings
in~\eqref{eq:three-cones} consequently has determinant $\pm1$.
Lemma~\ref{lem:fan} identifies the cone lattices with those generated by
the displayed rays, so these determinants are exactly the required
indices. The other two pairings collapse to dimension two.
\end{proof}

\section{Separating contributions and noncoordinate edges}

\begin{lemma}[Affine-span separation]\label{lem:separation}
Let $B=\sum_{j=1}^r c_j[P_j]\in\Ch_d(V)$, and fix an affine
$d$-plane $P$. The function
\[
 b_P=\sum_{j:\,\aff(P_j)=P}c_j\ind_{P_j}
\]
is determined almost everywhere by the chain $B$, independently of
its representation. Moreover, $B=0$ if and only if $b_P=0$ almost
everywhere for every affine span occurring in a representation.
Here measure on $P$ is Lebesgue measure on its direction space.
\end{lemma}

\begin{proof}
Every distinct affine $d$-plane $\aff(P_j)$ meets $P$ in dimension at
most $d-1$, if it meets it at all. Remove these finitely many intersections
and the boundaries of the polyhedra with affine span $P$. On what
remains, the displayed sum is exactly the coefficient of $B$ after
common subdivision. All removed sets have measure zero. Thus a zero
chain gives a zero function, and subtracting two representations proves
independence of the representation. Conversely, each full-dimensional
cell of a common subdivision has one of the finitely many affine spans.
Its coefficient is the almost-everywhere value of the corresponding
function. If all these functions vanish, all cell coefficients vanish.
\end{proof}

\begin{proposition}\label{prop:ordinary}
If $T(D)=0$, every nonzero edge of $D$ is parallel to some $e_i$.
\end{proposition}

\begin{proof}
Suppose that $e$ has nonzero weight $w_e$ and primitive direction $v$
not parallel to any $e_i$. Choose the partition and the space $U$ from
Lemma~\ref{lem:ordinary-linear}, and write $h=h_{ab}$.
Take an interior point $x_0$ of $e$ and a small neighborhood in which
$D$ is just a constant-weight open segment of $e$. This is possible
after common subdivision, since the chain has only finitely many cells.
Let $x$ vary over a smaller open segment about $x_0$, and set
\[
 z=\varphi(x)+\alpha E_{ab}+\beta E_{cd},
 \qquad \alpha,\beta>0,\quad \alpha+\beta<\varepsilon.
\]
For $\varepsilon>0$ small enough, these points form a nonempty open
subset of the affine three-plane $P=\varphi(x_0)+U$.

Any summand of~\eqref{eq:transform} with affine span $P$ must use a
cone contained in $U$. By Lemma~\ref{lem:ordinary-linear} this cone
is $\sigma_{ab|cd}$. Consider a competing representation
\[
 z=\varphi(y)+\alpha'E_{ab}+\beta'E_{cd},
 \qquad \alpha',\beta'\ge0,
\]
with $y$ on the support of $D$. Applying $\pi$ and using
\eqref{eq:identity} gives
\begin{equation}\label{eq:localization}
 \alpha'+\beta'=\alpha+\beta,\qquad
 y-x=(\alpha-\alpha')h,\qquad
 |\alpha-\alpha'|\le\alpha+\beta.
\end{equation}
By reducing $\varepsilon$ uniformly on the smaller segment, we force
$y$ into the neighborhood where the source consists only of $e$.
Since $v$ and $h$ are independent, the second identity then implies
$y=x$, $\alpha'=\alpha$, and $\beta'=\beta$.

Thus the sum of contributions with affine span $P$ has, on this open
set, the single coefficient $w_e m(e,\sigma_{ab|cd})\ne0$.
All other affine spans can be ignored by
Lemma~\ref{lem:separation}. This contradicts $T(D)=0$.
\end{proof}

\section{Thin bands and coordinate edges}

\subsection{A measure lemma}

Consider a finite signed rational one-chain $B$ in $\R^2$ whose edges
are horizontal or vertical. Its lattice-length measure is the locally
finite signed Radon measure, understood locally,
\begin{equation}\label{eq:source-measure}
 \mu_B=
 \sum_{(I,\eta)}w_{I,\eta}\,\ind_I(X)\,dX\,\delta_\eta(dY)
 +\sum_{(\xi,J)}w_{\xi,J}\,\delta_\xi(dX)\,\ind_J(Y)\,dY.
\end{equation}
The intervals may be bounded or unbounded. Endpoints have measure zero.
For $s\in\R$, let $\tau_s(X,Y)=(X+s,Y+s)$, and for $t>0$ put
\begin{align}
 C_t(X,Y)={}&
 \sum_{(I,\eta)}w_{I,\eta}
       \ind_{\{0<Y-\eta<t\}}\ind_I(X-Y+\eta)\notag\\
 &+\sum_{(\xi,J)}w_{\xi,J}
       \ind_{\{0<X-\xi<t\}}\ind_J(Y-X+\xi).
 \label{eq:central-density}
\end{align}

\begin{lemma}[Thin-band recovery]\label{lem:thin-band}
As locally finite signed Radon measures,
\begin{equation}\label{eq:band-measure}
 C_t\,dX\,dY=\int_0^t(\tau_s)_*\mu_B\,ds.
\end{equation}
As $t\downarrow0$,
\begin{equation}\label{eq:band-limits}
 C_t\longrightarrow0\quad\text{in }L^1_{\mathrm{loc}}(\R^2),
 \qquad
 t^{-1}C_t\,dX\,dY\longrightarrow\mu_B
\end{equation}
against all compactly supported continuous test functions.
The measure integral in~\eqref{eq:band-measure} is understood by
testing against compactly supported functions.
In particular, if $H\in L^1_{\mathrm{loc}}(\R^2)$ is independent of
$t$ and $H+C_t=0$ for almost every $(X,Y,t)$ with $t>0$, then $B=0$.
\end{lemma}

\begin{proof}
For a horizontal edge, the map $(x,s)\mapsto(x+s,\eta+s)$ has
Jacobian of absolute value one; for a vertical edge the same holds for
$(y,s)\mapsto(\xi+s,y+s)$. Change of variables gives
\eqref{eq:band-measure}, including the coefficients in
\eqref{eq:central-density}.

If $K\subset\R^2$ is compact and $0<t\le1$, then
\[
 \int_K|C_t|\,dX\,dY
 \le \int_0^t|\mu_B|(\tau_{-s}K)\,ds
 \le t\,|\mu_B|\bigl(K+[-1,0](1,1)\bigr).
\]
The last quantity is finite and tends to zero. This proves the first
limit even when the source has unbounded edges. For a compactly
supported continuous function $g$,
\[
 \frac1t\int g C_t\,dX\,dY
 =\int\left(\frac1t\int_0^t g(X+s,Y+s)\,ds\right)d\mu_B(X,Y).
\]
Uniform continuity on a fixed enlarged compact set shows that this
converges to $\int g\,d\mu_B$, proving the second limit.

For the last assertion, apply Fubini's theorem on a countable compact
exhaustion of $\R^2$. It gives a single full-measure set
$G\subset(0,1)$ such that $H+C_t=0$ almost everywhere on all of
$\R^2$ for each $t\in G$. Choose $t_n\in G$ decreasing to zero.
The first limit gives $H=0$ locally in $L^1$.
Then $C_{t_n}=0$, and the second limit gives $\mu_B=0$.
After common subdivision, an interior subsegment of any nonzero edge
can be isolated from all other edges. A nonnegative continuous test
function supported in this neighborhood and positive on the subsegment
has a nonzero integral against $\mu_B$, contradicting $\mu_B=0$.
Hence $B=0$.
\end{proof}

\subsection{The coordinate chart}

\begin{proposition}\label{prop:coordinate}
If every edge of $D$ is parallel to some $e_i$ and $T(D)=0$, then $D=0$.
\end{proposition}

\begin{proof}
Fix a pair $a,b$ with complement $c,d$, and an affine plane
$A=o+H_{ab}$ containing at least one source edge parallel to $e_a$
or $e_b$. Set $P=\varphi(o)+U_{ab}$. By
Lemma~\ref{lem:coordinate-linear}, the contributions with affine span
$P$ come exactly from the $e_a$- and $e_b$-parallel edges contained
in $A$, paired with the three cones in~\eqref{eq:three-cones}, and
with collapsed pairings omitted. Indeed, any contributing direction
$f_i$ and cone must lie in $U_{ab}$; its source offset then lies in
$A$ because $U_{ab}\cap\varphi(\NR)=\varphi(H_{ab})$.

Write points of $A$ as $o+Xe_a+Ye_b$, and let $B$ be the horizontal
and vertical chain given by these source edges. These coordinates
normalize lattice length, since $e_a,e_b$ form a basis of
$N\cap H_{ab}$. Use coordinates on $P$ given by
\begin{equation}\label{eq:chart}
 z=\varphi(o)+Xf_a+Yf_b+tE_{cd},\qquad t>0.
\end{equation}
All noncollapsed indices are one by
Lemma~\ref{lem:coordinate-linear}.

In the central cone, identity~\eqref{eq:identity} gives
\[
 rE_{ab}+sE_{cd}=r(f_a+f_b)+(r+s)E_{cd}.
\]
Thus, at thickness $t=r+s$, its planar displacement is $(r,r)$
with $0<r<t$. The total central contribution is precisely $C_t$
from~\eqref{eq:central-density}.

The two side cones give displacements $-r f_a$ and $-r f_b$,
while their $E_{cd}$ coefficient is $t$. A horizontal source edge
pairs nontrivially only with the $-f_b$ side, and a vertical source
edge only with the $-f_a$ side. Their total contribution is
\begin{equation}\label{eq:background}
 H(X,Y)=
 \sum_{(I,\eta)}w_{I,\eta}\ind_I(X)\ind_{\{Y<\eta\}}
 +\sum_{(\xi,J)}w_{\xi,J}\ind_{\{X<\xi\}}\ind_J(Y).
\end{equation}
This is a locally integrable function independent of $t$. It includes
all distant side-cone contributions, not just those near the point
being probed.

The weight function of the contributions with affine span $P$ is
therefore $K(X,Y,t)=H(X,Y)+C_t(X,Y)$, almost everywhere in the chart.
Since $T(D)=0$, Lemma~\ref{lem:separation} gives $K=0$ almost
everywhere. Lemma~\ref{lem:thin-band} implies $B=0$.

For a fixed pair there are only finitely many planes $A$ arising from
the source edges. Repeating the argument for these planes and for pairs
covering all four indices eliminates every edge of $D$.
\end{proof}

\begin{proof}[Proof of Theorems~\ref{thm:chains} and~\ref{thm:curves}]
If $T(D)=0$, Proposition~\ref{prop:ordinary} reduces the support to
coordinate directions, and Proposition~\ref{prop:coordinate} gives
$D=0$. This proves Theorem~\ref{thm:chains}. Given the hypotheses of
Theorem~\ref{thm:curves}, take the signed chain $D=\Sigma-\Sigma'$.
By Lemma~\ref{lem:well-defined}, its transform is
$Z_\Sigma-Z_{\Sigma'}=0$, so $D=0$.
\end{proof}

\begin{remark}
The thin-band argument uses only local finiteness of the length measure,
change of variables, and uniform continuity of test functions.
The conclusion does not require selecting an arbitrary slice of an
almost-everywhere identity: the sequence of valid slices in
Lemma~\ref{lem:thin-band} handles this explicitly.
\end{remark}

\section{A two-vertex example}

Let $v=e_0+e_1$, let $p\in N_{\Q}$, and let $L\in\Q_{>0}$.
Consider the chain consisting of the bounded edge $[p,p+Lv]$, the two
rays from $p$ in directions $-e_0,-e_1$, and the two rays from $p+Lv$
in directions $-e_2,-e_3$, all of weight $m\in\Z_{>0}$.
It is balanced at both vertices, since
\[
 v-e_0-e_1=0,\qquad -v-e_2-e_3=0.
\]
The parameters can be recovered without constructing the entire output.

\begin{proposition}\label{prop:example}
For this family the weighted Chow transform determines $p,L,m$.
\end{proposition}

\begin{proof}
Put $d=\varphi(v)$ and $U=\Span(d,E_{02},E_{13})$.
The split $02|13$ satisfies the hypothesis of
Lemma~\ref{lem:ordinary-linear}. Thus no $f_i$ belongs to $U$,
and the only incidence cone contained in $U$ is
$\sigma_{02|13}$. The output's three-dimensional part with direction
space $U$ consequently has support closure
\[
 S=\varphi(p)+[0,L]d+\pos(E_{02},E_{13}).
\]
This extraction is intrinsic to the weighted chain: contributions with
different direction spaces meet it only in lower dimension, and the
four unbounded source edges cannot contribute with direction space $U$.

The three vectors $d,E_{02},E_{13}$ are independent, so $S$ is an
affine image of $[0,L]\times\R_{\ge0}^2$. Its two vertices are
$\varphi(p)$ and $\varphi(p)+Ld$. The fixed positive direction $d$
orders these vertices, recovering $p$ and $L$. In the quotient
coordinates obtained by deleting $23$ after subtraction, the columns
$d,E_{02},E_{13}$ have a determinant-one minor in rows $03,02,13$.
The image weight is therefore $m$, recovering the last parameter.
\end{proof}

\begin{remark}
The finite hypothesis in the general theorem is used to separate affine
spans and to isolate source interiors. The theorem does not address
chains with infinitely many cells, higher-dimensional source varieties,
or additional projective boundary strata. The example gives an explicit
recovery procedure only for this family.
\end{remark}

\section*{Acknowledgments}

This mathematical paper was prepared with the assistance of artificial intelligence.  I
first acknowledge the textual data produced by humanity as a whole, which
provides the linguistic and cultural record from which mathematical expression
can be learned.  I then thank the accumulated scholarly achievements of
mathematicians throughout history.  I also thank the researchers whose work has
advanced and refined neural networks through successive generations.  Finally,
I thank the OpenAI team for the training and development of the models used in
the preparation of this paper.

\begin{thebibliography}{9}

\bibitem{AllermannHampeRau}
L.~Allermann, S.~Hampe, and J.~Rau,
\emph{On rational equivalence in tropical geometry},
\href{https://arxiv.org/abs/1408.1537v2}{arXiv:1408.1537v2}.

\bibitem{Fink}
A.~Fink,
\emph{Tropical cycles and Chow polytopes},
Beitr. Algebra Geom. \textbf{54} (2013), 13--40.
\href{https://arxiv.org/abs/1001.4784}{arXiv:1001.4784}.

\bibitem{Tripoli}
P.~Tripoli,
\emph{Tropical Chow hypersurfaces},
Int. Math. Res. Not. IMRN \textbf{2019}, no.~14, 4302--4324.
\href{https://doi.org/10.1093/imrn/rnx260}{doi:10.1093/imrn/rnx260}.
The numbered statements cited here are from
\href{https://arxiv.org/abs/1612.05192v2}{arXiv:1612.05192v2},
September 19, 2017.

\end{thebibliography}
\end{document}
